\Needspace{6\baselineskip}
\begin{center}\hypertarget{native-229-article166}{}\pdfbookmark[2]{166.}{article-166}166.\end{center}
For the first determination of a hitherto entirely unknown orbit from four observations, we must necessarily use positions that do not cover too large a heliocentric motion, since otherwise we would lack the means to form a first approximation readily. The method we are about to give, however, has so wide a scope that observations covering a heliocentric motion of $30$ or $40$ degrees can be used without hesitation, provided the distances from the Sun are not too unequal. Whenever we have a choice, it is helpful to take nearly equal time intervals between the first and second observations, the second and third, and the third and fourth. There is no need for undue concern in this respect either, as the example below will show. Its time intervals are $48$, $55$ and $59$ days, and its heliocentric motion exceeds $50^0$.

Furthermore, our solution requires the second and third observations to be complete, so the latitudes or declinations in the first and last observations are \SourceLeaf{230}neglected. We did point out above that, for the sake of greater precision, it is usually preferable to fit the elements to the first and last complete observations and to the longitudes or right ascensions of the intermediate ones. In the first determination of an orbit, however, we need not regret forgoing this advantage. Obtaining an approximation as readily as possible is of far greater importance. The loss falls mainly on the longitude of the node and the inclination of the orbit, has scarcely any perceptible effect on the remaining elements, and can easily be made good later.

For brevity we shall present the method by referring all positions to the ecliptic, and accordingly assume that four longitudes and two latitudes are given. Since our formulas also take the latitude of the Earth into account, however, they can readily be carried over to the case in which the equator is taken as the fundamental plane, provided right ascensions and declinations are substituted for longitudes and latitudes.

As for nutation, precession, parallax, and aberration, everything we explained in the preceding section applies here as well. Unless the approximate distances from the Earth are already known from another source, so that method I of Article 118 can be used for aberration, the observed positions will initially be corrected only for the aberration of the fixed stars. The times will then be corrected as soon as the distances can be determined approximately in the course of the calculation, as will become clearer below.
